\chapter{Análisis de Requisitos y Diseño}
\label{cap:diseno}
En este capítulo se describe la arquitectura del simulador desarrollado, detallando la jerarquía de clases y el flujo de la información desde la generación del mensaje hasta su decodificación.

\section{Especificación de Requisitos}
La especificación de requisitos establece el marco del trabajo y las capacidades que el simulador debe poseer para cumplir con los objetivos académicos y técnicos del proyecto.

\subsection{Requisitos Funcionales}
Los requisitos funcionales describen las tareas específicas que el software debe realizar:

\begin{itemize}

    \item \textbf{RF1. Aritmética en el Cuerpo de Galois:} El sistema debe implementar las operaciones de suma, multiplicación e inversión en el cuerpo $GF(2^m)$ mediante la tabla de búsqueda (\textit{log/antilog}).
    
    \item \textbf{RF2. Configuración de Parámetros:} El usuario podrá definir el grado del cuerpo ($m$) y la capacidad de corrección ($t$). El sistema calculará automáticamente los parámetros $(n,k)$ asociados y seleccionará el polinomio primitivo irreducible correspondiente al grado $m$ para construir el cuerpo.
    
    \item \textbf{RF3. Codificación sistemática:} El software deberá construir el polinomio generador $g(x)$ y generar palabras código sistemáticas en las que el mensaje original permanezca inalterado, seguido de los bits de paridad.
    
    \item \textbf{RF4. Modelado del Canal:} El sistema debe permitir la simulación de distintos modelos de canal, incluyendo el Canal Binario Simétrico (BSC) y el canal de Ruido Blanco Aditivo Gaussiano (AWGN) mediante modulación BPSK, realizando el mapeo de bits $\{0, 1\}$ a niveles de amplitud $\{1, -1\}$ para su transmisión por el canal.
    
    \item \textbf{RF5. Decodificación y Corrección:} El sistema deberá ser capaz de calcular los síndromes del mensaje recibido, generar el polinomio localizador de errores y corregir hasta $t$ errores de bit.

    \item \textbf{RF6. Procesamiento de Ficheros:} El simulador debe ser capaz de codificar, transmitir con ruido y decodificar bits extraídos de ficheros arbitrarios del sistema.
    
    \item \textbf{RF6. Registro de Métricas:} El sistema debe exportar a formato CSV los valores de BER \textit{(Bit Error Rate)}, BER sin codificar, CWER \textit{(Codeword Error Rate)} y el tiempo promedio de codificación y decodificación.
    
\end{itemize}

\subsection{Requisitos No Funcionales}
Los requisitos no funcionales definen las restricciones y atributos de calidad del sistema:
\begin{itemize}
    \item \textbf{RNF1. Eficiencia Temporal:} Dado el carácter estadístico de las simulaciones de Monte Carlo, las operaciones del simulador deben estar optimizadas en C++ para poder procesar grandes volúmenes de datos en un tiempo reducido.
    
    \item \textbf{RNF2. Modularidad:} El diseño debe separar el motor algebraico de la lógica de comunicaciones, permitiendo la reutilización de las clases de Cuerpos de Galois en otros esquemas de codificación. 
    
%    \item \textbf{RNF3. Portabilidad y Estándares:} 
\end{itemize}
    
\section{Diseño del Sistema}
Se ha diseñado el software bajo un enfoque modular, facilitando la validación de la lógica matemática y la lógica de simulación.
\subsection{Arquitectura}
La estructura se divide en tres niveles de abstracción:

\begin{enumerate}

    \item \textbf{Capa de Álgebra:} Contiene las abstracciones matemáticas de elementos del campo y polinomios. 
    
    \item \textbf{Capa de Aplicación (BCH):} Contiene el protocolo de codificación y decodificación, utilizando los servicios de la capa de álgebra.
    
    \item \textbf{Capa de Simulación:} Contiene el proceso de generación de ruido estadístico, la paralelización por hilos y la recolección de métricas de rendimiento.
    
\end{enumerate}

\subsection{Flujo de Información y Modelo de Simulación}
El flujo de información en el sistema se ha diseñado para seguir un modelo de \textit{pipeline} paralelo. A diCWERencia de un flujo secuencial convencional, el simulador distribuye la carga de trabajo en múltiples hilos de ejecución, donde cada uno gestiona de forma independiente el ciclo completo de vida de una palabra código.

En la figura \ref{fig:diagrama_flujo_paralelo} se ilustra este proceso. El flujo comienza con la lectura de la configuración global (\texttt{SimulationConfig}), que inicializa las tablas estáticas del cuerpo de Galois. Posteriormente, el orquestador lanza un pool de hilos donde cada uno realiza las etapas de generación, codificación, inyección de ruido y decodificación.

\begin{figure}[h]
    \centering
    \shorthandoff{>}
    \resizebox{\textwidth}{!}{
    \begin{tikzpicture}[
        node distance=1cm,
        block/.style={draw, rounded corners, minimum width=2.5cm, minimum height=1cm, align=center, font=\small},
        thread/.style={draw, dashed, gray, rounded corners, inner sep=10pt},
        arrow/.style={thick, ->, >=stealth}
    ]

    % Bloque inicial
    \node[block] (config) {Configuración\\y Tablas GF};

    % Representación de Paralelismo (Threads)
    \node[block, below=1.5cm of config] (enc1) {Codificación};
    \node[block, left=0.5cm of enc1] (enc0) {Codificación};
    \node[block, right=0.5cm of enc1] (enc2) {Codificación};
    
    \node[block, below=1.2cm of enc1] (canal1) {Canal\\(AWGN/BSC)};
    \node[block, left=0.5cm of canal1] (canal0) {Canal\\(AWGN/BSC)};
    \node[block, right=0.5cm of canal1] (canal2) {Canal\\(AWGN/BSC)};

    \node[block, below=1.2cm of canal1] (dec1) {Decodificación\\(B-M / Chien)};
    \node[block, left=0.5cm of dec1] (dec0) {Decodificación\\(B-M / Chien)};
    \node[block, right=0.5cm of dec1] (dec2) {Decodificación\\(B-M / Chien)};

    % Bloque final de reducción
    \node[block, below=1.5cm of dec1] (metrics) {Sincronización\\y Métricas (CSV)};

    % Caja que indica paralelismo
    \node[thread, fit=(enc0) (enc2) (dec0) (dec2), label=right:{\small Paralelismo OpenMP}] (pool) {};

    % Flechas
    \draw[arrow] (config) -- (pool.north);
    \draw[arrow] (enc0) -- (canal0); \draw[arrow] (canal0) -- (dec0);
    \draw[arrow] (enc1) -- (canal1); \draw[arrow] (canal1) -- (dec1);
    \draw[arrow] (enc2) -- (canal2); \draw[arrow] (canal2) -- (dec2);
    \draw[arrow] (pool.south) -- (metrics);

    \end{tikzpicture}
    }
    \shorthandon{>} 
    \caption{Flujo de información paralelo para simulación de Monte Carlo.}
    \label{fig:diagrama_flujo_paralelo}
\end{figure}

\subsection{Diseño de Datos y Estructuras}
\label{sub:diseno_datos}
Para garantizar el cumplimiento de los requisitos de rendimiento (RNF1) y modularidad (RNF2), se ha diseñado una jerarquía de clases para representar las abstracciones matemáticas en estructuras de datos eficientes en C++.
\subsubsection{Representación de la Información}
La unidad básica de información es el bit, que se representará de dos maneras distintas según el contexto:

\begin{itemize}
\item \textbf{Representación de elementos y coeficientes:} Para el manejo de coeficientes polinómicos y elementos del cuerpo, se utiliza \verb|uint16_t|. Esta elección se debe a que el simulador soporta hasta $m=15$ ($n=32.767$), rango que se ajusta perfectamente a la capacidad de 16 bits sin desperdicio de memoria. Al utilizar la mitad de espacio que un entero estándar de 32 bits, se mejora la localidad de datos y se aprovecha mejor la memoria caché del procesador, acelerando las iteraciones de Monte Carlo.

\item \textbf{Empaquetado de Bit (\textit{Bit-packing}):} En procesos como la división polinómica binaria del codificador, se utiliza la estructura \verb|PackedBit|. Esta permite procesar bloques de bits simultáneamente, reduciendo significativamente los ciclos de CPU.

\item \textbf{Tablas de búsqueda del Campo de Galois:} La clase \verb|GaloisField| utiliza miembros estáticos para almacenar \verb|log_table| y \verb|exp_table|. Al ser \verb|static|, las tablas se calculan una sola vez para un valor de $m$ dado y son compartidas por todas las instancias, optimizando el uso de la memoria caché.

\item \textbf{Gestión de coeficientes en \texttt{Polynomial}:} Los polinomios se representan mediante un \verb|std::vector<uint16_t>| donde el índice coincide con el grado del término. Se incluye un método \verb|trim()| que elimina los ceros de mayor grado para asegurar que el tamaño del vector sea siempre el mínimo necesario.

\item \textbf{Cálculo de síndromes:} En la clase \verb|BCH_Codec|, el cálculo de síndromes se apoya en un \verb|std::vector| de tamaño $2t+1$. Se aprovecha la propiedad de los cuerpos binarios para calcular los síndromes de índice par como $S_i = (S_{i/2})^2$, ahorrando la mitad de las evaluaciones polinómicas.

\item \textbf{Estructuras de decodificación:} Durante el algoritmo de Berlekamp-Massey y la búsqueda de Chien, se emplean vectores auxiliares como \verb|reg| y \verb|factors| para almacenar los coeficientes del polinomio localizador de errores y sus raíces inversas, permitiendo una corrección de errores en tiempo lineal respecto a $n$.
    
\item \textbf{Generación de ruido aleatorio:} La clase \verb|Channel| encapsula un motor de generación \verb|std::mt19937| (Mersenne Twister). El uso de este generador, junto con \verb|std::random_device| para la semilla, asegura que las muestras de ruido para los modelos BSC y AWGN sean estadísticamente independientes en cada iteración de Monte Carlo.

\end{itemize}

\subsection{Diseño Orientado a Objetos}
A partir de los requisitos y estructuras de datos definidas, se ha diseñado una arquitectura orientada a objetos. En la figura \ref{fig:uml_clases} se detalla la jerarquía y los atributos de las clases implementadas. A continuación se describen las relaciones entre las clases principales:  
\begin{itemize}
    \item \textbf{Relación de Composición:} La clase \verb|BCH_Codec| mantiene una instancia de \texttt{GaloisField} y una de \verb|Polynomial| (el polinomio generador). Esta jerarquía asegura que el código BCH siempre opere bajo las reglas aritméticas del cuerpo definido.
    \item \textbf{Independencia de Canal:} La clase \verb|Channel| se ha diseñado de manera desacoplada del \verb|BCH_Codec|, por lo que simplemente recibe un vector de bits y devuelve una versión ruidosa.
\end{itemize}

\begin{figure}[h]
    \centering
    \includegraphics[width=0.85\textwidth]{imagenes/UML.png}
    \caption{Diagrama de clases UML del simulador BCH.}
    \label{fig:uml_clases}
\end{figure}

Debido a la densidad de la jerarquía de clases y para facilitar su análisis técnico en esta memoria, la arquitectura se ha descompuesto en tres subsistemas funcionales. Esta subdivisión permite observar con mayor claridad las relaciones de dependencia, composición y la especialización de cada componente dentro del simulador.

\subsubsection{Subsistema Matemático y de Optimización}
Este primer bloque constituye la base del simulador. Engloba las clases responsables de la aritmética en el cuerpo de Galois y la manipulación de polinomios. 

Como se observa en la figura \ref{fig:uml_math}, la clase \texttt{Polynomial} mantiene una relación de composición con \texttt{GaloisField}, asegurando que todas las operaciones polinómicas respeten las reglas del campo definido. Además, se incluye la clase \texttt{PackedBit} como estructura de soporte para optimizar las operaciones de desplazamiento y XOR a nivel de bit, críticas para el rendimiento del sistema.

\begin{figure}[h]
    \centering
    \includegraphics[width=0.75\textwidth]{Plantilla_TFG_latex/imagenes/uml_mat.png} 
    \caption{Diagrama de clases: Subsistema Matemático.}
    \label{fig:uml_math}
\end{figure}

\subsubsection{Subsistema Core Codec BCH}
El segundo nivel de abstracción corresponde a la implementación específica del código BCH. La clase \texttt{BCH\_Codec} actúa como una fachada (\textit{Facade}) que expone los métodos de codificación y decodificación, ocultando la complejidad de los algoritmos de Berlekamp-Massey y la búsqueda de Chien.

La figura \ref{fig:uml_codec} detalla cómo esta clase orquesta el uso de los polinomios y el campo de Galois para calcular el polinomio generador durante la fase de inicialización y corregir los errores en la fase de recepción.

\begin{figure}[h]
    \centering
    \includegraphics[width=0.6\textwidth]{Plantilla_TFG_latex/imagenes/uml_codec.png}
    \caption{Diagrama de clases: Núcleo del Códec BCH.}
    \label{fig:uml_codec}
\end{figure}

\subsubsection{Subsistema de Orquestación y Simulación}
Finalmente, el nivel superior gestiona el entorno de ejecución y la recolección de métricas. La clase \texttt{BCH\_Simulator} es la encargada de inyectar la configuración (\texttt{SimulationConfig}) y gestionar el canal de comunicaciones.

Como muestra la figura \ref{fig:uml_sim}, este subsistema desacopla el modelo de ruido (\texttt{Channel}) del algoritmo de corrección, permitiendo que las métricas de rendimiento se recojan de forma estructurada en objetos \texttt{PointResults} para su posterior exportación a CSV.

\begin{figure}[h]
    \centering
    \includegraphics[width=0.85\textwidth]{Plantilla_TFG_latex/imagenes/uml_sim.png}
    \caption{Diagrama de clases: Entorno de simulación y orquestación.}
    \label{fig:uml_sim}
\end{figure}

\section{Metodología de Simulación}
El análisis del rendimiento se basa en el método de \textbf{Monte Carlo}, un proceso iterativo que permite estimar Probabilidades mediante la repetición de experimentos aleatorios. Para garantizar la convergencia estadística de los resultados, la simulación de cada punto de relación señal-ruido continúa hasta que se contabiliza un mínimo de 100 tramas erróneas o se alcanza el límite máximo de iteraciones configurado.

\subsection{Diseño del Bucle de Simulación}
El simulador se ha diseñado para ejecutar el siguiente flujo iterativo hasta alcanzar un nivel de confianza estadística suficiente:
\begin{enumerate}
    \item \textbf{Generación:} Creación de un mensaje aleatorio de $k$ bits.
    \item \textbf{Codificación:} Transformación del mensaje en una palabra código de $n$ bits mediante \verb|BCH_Codec|.
    \item \textbf{Transmisión:} Inyección de ruido en el \verb|Channel| (BSC o AWGN) según el nivel de ruido configurado.
    \item \textbf{Decodificación:} Intento de recuperación del mensaje original.
    \item \textbf{Comparación:} Contabilización de errores de bit (BER) y errores de trama (CWER) comparando el mensaje original con el decodificado.
\end{enumerate}

\subsection{Parámetros de Control}
Para el diseño de las pruebas, se han definido variables de control que permiten analizar la escalabilidad del sistema: el tamaño del campo ($m$), la capacidad de corrección ($t$) y el número de iteraciones ($N$), el cual determina la precisión de los resultados estadísticos obtenidos.

\section{Metodología de Simulación}
El análisis del rendimiento se basa en el método de \textbf{Monte Carlo}, un proceso iterativo que permite estimar Probabilidades mediante la repetición de experimentos aleatorios. Para garantizar la convergencia estadística de los resultados, la simulación de cada punto de relación señal-ruido continúa hasta que se contabiliza un mínimo de 100 tramas erróneas o se alcanza el límite máximo de iteraciones configurado.

\subsection{Diseño del Bucle de Simulación}
El simulador se ha diseñado para ejecutar el siguiente flujo iterativo hasta alcanzar un nivel de confianza estadística suficiente:
\begin{enumerate}
    \item \textbf{Generación:} Creación de un mensaje aleatorio de $k$ bits.
    \item \textbf{Codificación:} Transformación del mensaje en una palabra código de $n$ bits mediante \verb|BCH_Codec|.
    \item \textbf{Transmisión:} Inyección de ruido en el \verb|Channel| (BSC o AWGN) según el nivel de ruido configurado.
    \item \textbf{Decodificación:} Intento de recuperación del mensaje original.
    \item \textbf{Comparación:} Contabilización de errores de bit (BER) y errores de trama (CWER) comparando el mensaje original con el decodificado.
\end{enumerate}

\subsection{Parámetros de Control}
Para el diseño de las pruebas, se han definido variables de control que permiten analizar la escalabilidad del sistema: el tamaño del campo ($m$), la capacidad de corrección ($t$) y el número de iteraciones ($N$), el cual determina la precisión de los resultados estadísticos obtenidos.
\section{Planificación}
\begin{itemize}
    \item \textbf{Sprint 1 - Motor Algebraico:} Implementación de $GF(2^m)$ y las aritméticas básicas de los polinomios en dicho cuerpo.
    \item \textbf{Sprint 2 - Lógica BCH:} Desarrollo de algoritmo de codificación y decodificación (Berlekamp-Massey).
    \item \textbf{Sprint 3 - Entorno de Simulación:} Creación de las clases del canal y exportación de los datos.
    \item \textbf{Sprint 4 - Análisis y Visualización:} Desarrollo de scripts en Python para el post-procesamiento y generación de curvas de error.
\end{itemize}
